\section{The Golem Myth}
\label{sec:introduction_the_golem_myth}

A creation that surpasses its creator.

Long before modern technology, humanity had already questioned the possibility of losing control over its creations.
One of the earliest expressions of this idea can be found in the story of the Golem.
In Hebrew tradition, the term Golem refers to an artificial creature created by humans \parencite{nocks1998golem}.
These creatures are made of clay and brought to life using instructions inscribed on their bodies, as illustrated in Figure~\ref{fig:golem_creation}.
The earliest accounts describe the Golem as a protective and helpful figure.

\begin{figure}[t]
    \centering
    \includegraphics[width=0.75\textwidth]{figures/introduction/golem_creation.png}
    \caption[The creation of a Golem.]{The creation of a Golem \parencite{ales1908venec}.}
    \label{fig:golem_creation}
\end{figure}


This vision becomes more ambiguous in the seventeenth century, as illustrated by the legend of the Golem of Chelm.
In this story, the initially small Golem gradually grows in size and strength until it can no longer be safely controlled.
The creation intended to protect others therefore becomes a potential threat to them.
The only remaining solution is to destroy the creature in order to protect the community.

The Golem has inspired similar narratives over time, including the modern notion of the robot.
The concept of the robot as an artificial servant was introduced in the play \textit{Rossum's Universal Robots}, written in 1920 \parencite{capek1920rur}.
In the play, a scientist creates artificial beings designed to perform labor.
In this first version, robots are organic creatures forged through advances in chemistry.
Conceived as servants, these artificial beings turn against humanity and destroy civilization.

The connection between the Golem and the robot was explicitly made by the author, as illustrated by his statement that \enquote{the Robot is the Golem made flesh through mass production} \parencite{comrada1995golem}.
Despite their destructive role in the play, the word robot is now used to designate artificial beings created to serve humans.
This shows that, from its origin, the word robot has embodied the hopes and fears surrounding technological progress.

Today, this tension is particularly evident in the field of artificial intelligence.
With the rise of computer technology, organic robots were gradually replaced in the collective imagination by mechanical ones, but the underlying fears remained.
These concerns are even stronger as progress in the field has continued to accelerate.
In these times of technological revolution, no one can predict how this development will unfold.
What is certain, however, is that this field offers an almost literal adaptation of the Golem myth.

As artificial intelligence systems become increasingly capable, they are also granted increasing autonomy.
They no longer simply execute predefined instructions, but are now trusted to decide for themselves how to achieve the goals they are given.
This autonomy is what makes them powerful, but it also makes their behavior harder to control.
Understanding the behavior of such systems is the objective of this thesis.
This work therefore places itself within a long-standing line of thought, one that can be summarized by a single question:
\begin{center}
\textit{Can we remain in control of what we create?}
\end{center}